
\documentclass{article}
\usepackage{colortbl}
\usepackage{makecell}
\usepackage{multirow}
\usepackage{supertabular}

\begin{document}

\newcounter{utterance}

\centering \large Interaction Transcript for game `ifeval', experiment `ifeval', episode 125 with reflexion{-}llama{-}3.1{-}8b{-}instruct.
\vspace{24pt}

{ \footnotesize  \setcounter{utterance}{1}
\setlength{\tabcolsep}{0pt}
\begin{supertabular}{c@{$\;$}|p{.15\linewidth}@{}p{.15\linewidth}p{.15\linewidth}p{.15\linewidth}p{.15\linewidth}p{.15\linewidth}}
    \# & \multicolumn{2}{c}{Player} && \multicolumn{2}{c}{Game Master} \\
    \hline

    \theutterance \stepcounter{utterance}  
    & & & \multicolumn{4}{p{0.6\linewidth}}{
        \cellcolor[rgb]{0.9,0.9,0.9}{
            \makecell[{{p{\linewidth}}}]{
                \texttt{\tiny{[P1$\langle$GM]}}
                \texttt{You have attempted this same task before and did not succeed. These are the plans you wrote after those attempts.} \\
\texttt{Trial \#1: New plan: I failed to provide two distinct and compelling product descriptions. In my first description, I focused on the practical benefits of the edible paper, such as reducing waste and adding creativity to meal presentation. However, I didn't emphasize the unique selling point of the product being edible. In my second description, I highlighted the sensory experience of using the edible paper, but I didn't provide enough details about the product's features and benefits.} \\
\\ 
\texttt{Next time, I will focus on creating two descriptions that showcase the product's unique aspects and appeal to different target audiences. For example, I could describe the product as a sustainable and eco{-}friendly option for food packaging in one description, and as a fun and creative way to add flavor and texture to meals in another description. I will also make sure to include specific details about the product's features, such as its taste, texture, and durability, to make the descriptions more engaging and convincing.} \\
\\ 
\\ 
\texttt{Generate two alternative product descriptions: The product is a new type of paper that can be used to wrap food, and is edible.} \\
\texttt{First repeat the prompt above without change, then give your answer. Please do not say any word before repeating the prompt above.} \\
            }
        }
    }
     \\ \\

    \theutterance \stepcounter{utterance}  
    & \multicolumn{4}{p{0.6\linewidth}}{
        \cellcolor[rgb]{0.561,0.737,0.561}{
            \makecell[{{p{\linewidth}}}]{
                \texttt{\tiny{[P1$\rangle$GM]}}
                \texttt{You have attempted this same task before and did not succeed. These are the plans you wrote after those attempts.} \\
\\ 
\texttt{Trial \#1: New plan: I failed to provide two distinct and compelling product descriptions. In my first description, I focused on the practical benefits of the edible paper, such as reducing waste and adding creativity to meal presentation. However, I didn't emphasize the unique selling point of the product being edible. In my second description, I highlighted the sensory experience of using the edible paper, but I didn't provide enough details about the product's features and benefits.} \\
\\ 
\texttt{Next time, I will focus on creating two descriptions that showcase the product's unique aspects and appeal to different target audiences. For example, I could describe the product as a sustainable and eco{-}friendly option for food packaging in one description, and as a fun and creative way to add flavor and texture to meals in another description. I will also make sure to include specific details about the product's features, such as its taste, texture, and durability, to make the descriptions more engaging and convincing.} \\
\\ 
\texttt{Product Description 1: } \\
\texttt{Introducing EdiWrap, the revolutionary edible paper that's changing the way we think about food packaging. Made from natural ingredients and biodegradable materials, EdiWrap is not only good for the planet but also for your taste buds. With its delicate, slightly sweet flavor and tender texture, EdiWrap adds a touch of elegance to any meal. Whether you're wrapping sushi, sandwiches, or snacks, EdiWrap is the perfect sustainable alternative to traditional paper or plastic wraps. Its durability ensures that your food stays fresh and secure, while its edibility means you can enjoy the packaging itself. Say goodbye to waste and hello to a more sustainable future with EdiWrap.} \\
\\ 
\texttt{Product Description 2: } \\
\texttt{Get ready to take your mealtime to the next level with EdiWrap, the edible paper that's as fun as it is functional. Imagine wrapping your favorite snacks or meals in a delicate, edible sheet that's almost too pretty to eat. Almost. EdiWrap's unique flavor and texture will add a new dimension to your culinary creations, and its versatility makes it perfect for everything from appetizers to desserts. Whether you're a foodie, a chef, or just a lover of all things edible, EdiWrap is the ultimate creative tool. So why settle for boring, disposable packaging when you can have a delicious, Instagram{-}worthy wrap that's as much a part of the meal as the food itself?} \\
            }
        }
    }
    & & \\ \\

    \theutterance \stepcounter{utterance}  
    & & & \multicolumn{2}{p{0.3\linewidth}}{
        \cellcolor[rgb]{0.9,0.9,0.9}{
            \makecell[{{p{\linewidth}}}]{
                \texttt{\tiny{[GM$|$GM]}}
                {'combination:repeat_prompt': {'prompt_to_repeat': 'Generate two alternative product descriptions: The product is a new type of paper that can be used to wrap food, and is edible.'}}
            }
        }
    }
    & & \\ \\

    \theutterance \stepcounter{utterance}  
    & & & \multicolumn{2}{p{0.3\linewidth}}{
        \cellcolor[rgb]{0.9,0.9,0.9}{
            \makecell[{{p{\linewidth}}}]{
                \texttt{\tiny{[GM$|$GM]}}
                \texttt{game\_result = LOSE} \\
            }
        }
    }
    & & \\ \\

    \theutterance \stepcounter{utterance}  
    & & \multicolumn{4}{p{0.6\linewidth}}{
        \cellcolor[rgb]{0.95,0.95,0.95}{
            \makecell[{{p{\linewidth}}}]{
                \texttt{\tiny{[P1$\leftrightarrow$P]}}
                \texttt{Post{-}episode Reflexion (not sent to GM):} \\
\texttt{New plan: I failed to provide two distinct and compelling product descriptions that effectively highlighted the unique aspects of the edible paper. In both descriptions, I focused on the benefits of the product, but I didn't create a clear contrast between the two descriptions. } \\
\\ 
\texttt{Next time, I will create two descriptions that not only showcase the product's unique aspects but also appeal to different target audiences. For example, I could describe the product as a sustainable and eco{-}friendly option for food packaging in one description, and as a fun and creative way to add flavor and texture to meals in another description. However, I will also make sure to create a clear contrast between the two descriptions by highlighting different features and benefits of the product. } \\
\\ 
\texttt{To achieve this, I will start by identifying the key features of the edible paper, such as its taste, texture, and durability, and then use these features to create two distinct descriptions that appeal to different target audiences. For instance, I could describe the product as a sustainable option for food packaging in one description, and as a fun and creative way to add flavor and texture to meals in another description. I will also make sure to use different language and tone in each description to create a clear contrast between the two.} \\
            }
        }
    }
    & \\ \\

\end{supertabular}
}

\end{document}
